% GENERATED FILE. Edit canonical lessons, then run npm run generate:books.
% canonical-source-hash: fnv1a64:e1244a14e1adf53a
% canonical-lessons: ML-C68-sixth, ML-C68-seventh, ML-C68-eighth, ML-C68-ninth, ML-C68-tenth, ML-C68-eleventh, ML-C68-at-first

\chapter{All the Way Up, and First in Time}
\label{ch:ml-ordinals-to-eleventh}
\hlchaptermodality{68}

\begin{chapteropening}
\textbf{By the end of this chapter:} \emph{I can say sixth to eleventh in Malayalam, and put an action first in time as well as a thing first in a line.}

Order a thing in a line with onnaam and an action in time with aadyam, and say which question each one answers.
\end{chapteropening}

\section[āṟāṁ]{\ml{ആറാം} (āṟāṁ) — sixth}

\label{lesson:ML-C68-sixth}



\begin{quote}
The line goes on past five, so the chapter opens by retrieving where it stopped.

\begin{itemize}
\raggedright
  \item \emph{Recall:} say \emph{añcāṁ} — \textbf{R1}
  \item \emph{Recall:} say \emph{onnāṁ}, and say what you did to \emph{onnŭ} to build it — \textbf{R2}
  \item \emph{Recall:} say \emph{pinne} — \textbf{R3}
  \item \emph{Recall:} say \emph{vātil} — \textbf{R4}
\end{itemize}
\end{quote}

\subsection*{You'll want to know: \ml{ആറാം}}
\textbf{\ml{ആറാം}} (\emph{āṟāṁ}) — \textquotedblleft{}\textbf{sixth}\textquotedblright{}. \textbf{\ml{ആറ്}} drops its \textbf{ŭ} and takes \textbf{\ml{ആം}}, with nothing else to notice.

\textbf{\ml{ആറാം} \ml{ദിവസം}} (\emph{āṟāṁ divasam}), the sixth day.

That \textquotedblleft{}nothing else to notice\textquotedblright{} is the point of the whole chapter. Six to ten were a separate lesson from one to five when you learned to count, because the words themselves had to be learned. For \textbf{ordering} they are not a separate anything.

\subsection*{Your turn}
Say these aloud:

\begin{itemize}
\raggedright
  \item \emph{āṟŭ}, then \emph{āṟāṁ}
  \item \emph{añcāṁ, āṟāṁ}
  \item \emph{āṟāṁ divasam} — the sixth day
\end{itemize}

\begin{tcolorbox}[breakable,colback=teal!4,colframe=teal!35!black,title={Before you move on}]
Say \textquotedblleft{}the sixth day\textquotedblright{}. (\textbf{\ml{ആറാം} \ml{ദിവസം}}.) Did the ending change above five? (\textbf{No}.)
\end{tcolorbox}
\section[ēḻāṁ]{\ml{ഏഴാം} (ēḻāṁ) — seventh}

\label{lesson:ML-C68-seventh}



\begin{quote}
\begin{itemize}
\raggedright
  \item \emph{Recall:} say \emph{āṟāṁ} — \textbf{R1}
  \item \emph{Recall:} say \emph{raṇṭāṁ}, and say which word Malayalam did NOT use for \textquotedblleft{}first\textquotedblright{} — \textbf{R2}
  \item \emph{Recall:} say \emph{uṭane} — \textbf{R3}
  \item \emph{Recall:} say \emph{kasēra} — \textbf{R4}
\end{itemize}
\end{quote}

\subsection*{You'll want to know: \ml{ഏഴാം}}
\textbf{\ml{ഏഴാം}} (\emph{ēḻāṁ}) — \textquotedblleft{}\textbf{seventh}\textquotedblright{}.

\textbf{\ml{ഏഴ്}} carries the curled \textbf{ḻ}, the sound in the name \emph{Malayāḷam} itself and the one this language shares with Tamil while Kannada and Telugu hardened it away. Adding the ending does not soften it: \textbf{\ml{ഏഴാം}}, with the tongue still curled back.

\textbf{\ml{ഏഴാം} \ml{ദിവസം}} (\emph{ēḻāṁ divasam}), the seventh day.

One honest warning about the page. This word begins with \textbf{\ml{ഏ}}, the independent long \textbf{ē}, and that is a letter this book has taught you to \emph{recognise} but not yet to \emph{write} — no stroke order for it has been sourced, so none is offered. Read it; do not copy it yet. That is why this lesson is headed by its romanization rather than by the Malayalam word, the way the counting lessons were.

\subsection*{Your turn}
\begin{itemize}
\raggedright
  \item \emph{Say it:} \emph{ēḻŭ}, then \emph{ēḻāṁ} — keep the ḻ curled both times
  \item \emph{Say it:} \emph{āṟāṁ, ēḻāṁ}
  \item \emph{Look:} at \ml{ഏഴാം} and find the \ml{ഴ}
\end{itemize}

\begin{tcolorbox}[breakable,colback=teal!4,colframe=teal!35!black,title={Before you move on}]
Say \textquotedblleft{}the seventh day\textquotedblright{}. (\textbf{\ml{ഏഴാം} \ml{ദിവസം}}.) What happens to the \ml{ഴ} when the ending goes on? (\textbf{Nothing}.)
\end{tcolorbox}
\section[eṭṭāṁ]{\ml{എട്ടാം} (eṭṭāṁ) — eighth}

\label{lesson:ML-C68-eighth}



\begin{quote}
\begin{itemize}
\raggedright
  \item \emph{Recall:} say \emph{ēḻāṁ} — \textbf{R1}
  \item \emph{Recall:} say \emph{mūnnāṁ} — \textbf{R2}
  \item \emph{Recall:} say \emph{cilappōḷ} — \textbf{R3}
  \item \emph{Recall:} say \emph{kōlaṁ} — \textbf{R4}
\end{itemize}
\end{quote}

\subsection*{You'll want to know: \ml{എട്ടാം}}
\textbf{\ml{എട്ടാം}} (\emph{eṭṭāṁ}) — \textquotedblleft{}\textbf{eighth}\textquotedblright{}.

\textbf{\ml{എട്ട്}} has a doubled retroflex \textbf{ṭṭ}, written as \ml{ട} stacked under \ml{ട} with the virama between them. The ending lands after all of that, so the stack is untouched: \textbf{\ml{എട്ടാം}}.

\textbf{\ml{എട്ടാം} \ml{ദിവസം}} (\emph{eṭṭāṁ divasam}), the eighth day.

\subsection*{Your turn}
\begin{itemize}
\raggedright
  \item \emph{Say it:} \emph{eṭṭŭ}, then \emph{eṭṭāṁ}
  \item \emph{Say it:} \emph{ēḻāṁ, eṭṭāṁ}
  \item \emph{Look:} at \ml{എട്ടാം} and find the doubled \ml{ട്ട} still in place
\end{itemize}

\begin{tcolorbox}[breakable,colback=teal!4,colframe=teal!35!black,title={Before you move on}]
Say \textquotedblleft{}the eighth day\textquotedblright{}. (\textbf{\ml{എട്ടാം} \ml{ദിവസം}}.) Did the doubled \ml{ട്ട} change? (\textbf{No}.)
\end{tcolorbox}
\section[ompathāṁ]{\ml{ഒമ്പതാം} (ompathāṁ) — ninth}

\label{lesson:ML-C68-ninth}



\begin{quote}
\begin{itemize}
\raggedright
  \item \emph{Recall:} say \emph{eṭṭāṁ} — \textbf{R1}
  \item \emph{Recall:} say \emph{nālāṁ} — \textbf{R2}
  \item \emph{Recall:} say \emph{mātraṁ} — \textbf{R3}
  \item \emph{Recall:} say \emph{pūvŭ} — \textbf{R4}
\end{itemize}
\end{quote}

\subsection*{You'll want to know: \ml{ഒമ്പതാം}}
\textbf{\ml{ഒമ്പതാം}} (\emph{ompathāṁ}) — \textquotedblleft{}\textbf{ninth}\textquotedblright{}.

\textbf{\ml{ഒമ്പത്}} was already a built word when you learned it: an old word for \textbf{one} plus \emph{pathŭ}, \textbf{ten} — one short of ten, with the \emph{n} turning to \emph{m} before the \emph{p}. The ordinal ending goes on top of that, and none of the machinery underneath reacts to it.

\textbf{\ml{ഒമ്പതാം} \ml{ദിവസം}} (\emph{ompathāṁ divasam}), the ninth day.

\subsection*{Your turn}
Say these aloud:

\begin{itemize}
\raggedright
  \item \emph{ompathŭ}, then \emph{ompathāṁ}
  \item \emph{eṭṭāṁ, ompathāṁ}
  \item \emph{on-pathŭ} slowly, then let it settle into \emph{ompathāṁ}
\end{itemize}

\begin{tcolorbox}[breakable,colback=teal!4,colframe=teal!35!black,title={Before you move on}]
Say \textquotedblleft{}the ninth day\textquotedblright{}. (\textbf{\ml{ഒമ്പതാം} \ml{ദിവസം}}.) What two words are inside \ml{ഒമ്പത്}? (An old \textbf{one}, and \emph{pathŭ}, \textbf{ten}.)
\end{tcolorbox}
\section[pathāṁ]{\ml{പത്താം} (pathāṁ) — tenth}

\label{lesson:ML-C68-tenth}



\begin{quote}
\begin{itemize}
\raggedright
  \item \emph{Recall:} say \emph{ompathāṁ} — \textbf{R1}
  \item \emph{Recall:} say \emph{añcāṁ} — \textbf{R2}
  \item \emph{Recall:} say \emph{ennāl} — \textbf{R3}
  \item \emph{Recall:} say \emph{māla} — \textbf{R4}
\end{itemize}
\end{quote}

\subsection*{You'll want to know: \ml{പത്താം}}
\textbf{\ml{പത്താം}} (\emph{pathāṁ}) — \textquotedblleft{}\textbf{tenth}\textquotedblright{}. \textbf{\ml{പത്ത്}} loses its \textbf{ŭ} and takes \textbf{\ml{ആം}}, and the ten are done.

\textbf{\ml{പത്താം} \ml{ദിവസം}} (\emph{pathāṁ divasam}), the tenth day.

\subsection*{Your turn}
Say these aloud:

\begin{itemize}
\raggedright
  \item \emph{pathŭ}, then \emph{pathāṁ}
  \item \emph{ompathāṁ, pathāṁ}
  \item all ten in order, from \emph{onnāṁ}
\end{itemize}

\begin{tcolorbox}[breakable,colback=teal!4,colframe=teal!35!black,title={Before you move on}]
Say \textquotedblleft{}the tenth day\textquotedblright{}. (\textbf{\ml{പത്താം} \ml{ദിവസം}}.) Count the ordinals from one to ten. How many of the ten were words you had to learn? (\textbf{None}.)
\end{tcolorbox}
\section[pathinonnāṁ]{\ml{പതിനൊന്നാം} (pathinonnāṁ) — eleventh, and the edge that is not there}

\label{lesson:ML-C68-eleventh}



\begin{quote}
\begin{itemize}
\raggedright
  \item \emph{Recall:} say \emph{pathāṁ} — \textbf{R1}
  \item \emph{Recall:} say \emph{āṟāṁ} — \textbf{R2}
  \item \emph{Recall:} say \emph{apēkṣa} — \textbf{R3}
  \item \emph{Recall:} say \emph{ākāśaṁ} — \textbf{R4}
\end{itemize}
\end{quote}

\subsection*{You'll want to know: \ml{പതിനൊന്നാം}}
\textbf{\ml{പതിനൊന്നാം}} (\emph{pathinonnāṁ}) — \textquotedblleft{}\textbf{eleventh}\textquotedblright{}.

\textbf{\ml{പതിനൊന്ന്}} was itself a built word: \emph{pathin-}, the joining shape of ten, with \emph{onnŭ} welded on. Now the ordinal ending goes on top of \emph{that}, and the result is a word of three parts, not one of which had to be learned as a whole.

\textbf{\ml{പതിനൊന്നാം} \ml{ദിവസം}} (\emph{pathinonnāṁ divasam}), the eleventh day.

\begin{grammarlens}[title={Grammar Lens: the ending is a fact about numbers, not about ten}]
Most books stop an ordinal chapter at \emph{tenth}, and a learner reasonably concludes that those ten words were the set. In Malayalam there is no set. \textbf{\ml{ആം}} attaches to a \textbf{number}, so every number you learn from here on arrives with its ordinal already made: \textbf{\ml{ഇരുപത്}} gives \textbf{\ml{ഇരുപതാം}} (\emph{irupathāṁ}), twentieth, without a lesson.

That is why this book taught the ending in the very first lesson of the previous chapter rather than after ten examples. The examples were never the content.
\end{grammarlens}

\subsection*{Your turn}
Say these aloud:

\begin{itemize}
\raggedright
  \item \emph{pathinonnŭ}, then \emph{pathinonnāṁ}
  \item \emph{pathāṁ, pathinonnāṁ}
  \item take \emph{irupathŭ} and make its ordinal yourself
\end{itemize}

\begin{tcolorbox}[breakable,colback=teal!4,colframe=teal!35!black,title={Before you move on}]
Say \textquotedblleft{}the eleventh day\textquotedblright{}. (\textbf{\ml{പതിനൊന്നാം} \ml{ദിവസം}}.) Do you need a lesson for \emph{twentieth}? (\textbf{No} — the ending goes on any number.)
\end{tcolorbox}
\section[Practice]{\ml{ആദ്യം} (ādyaṁ) — at first, and the two kinds of order}

\label{lesson:ML-C68-at-first}



\begin{quote}
\begin{itemize}
\raggedright
  \item \emph{Recall:} say \emph{pathinonnāṁ}, and say why \emph{irupathāṁ} needs no lesson — \textbf{R1}
  \item \emph{Recall:} say \emph{ēḻāṁ} — \textbf{R2}
  \item \emph{Recall:} say \emph{anuvādaṁ} — \textbf{R3}
  \item \emph{Recall:} say \emph{sūryan} — \textbf{R4}
\end{itemize}
\end{quote}

\subsection*{You'll want to know: \ml{ആദ്യം}}
\textbf{\ml{ആദ്യം}} (\emph{ādyaṁ}) — \textquotedblleft{}\textbf{at first}\textquotedblright{}, \textquotedblleft{}\textbf{first of all}\textquotedblright{}.

This is not another ordinal. \textbf{\ml{ഒന്നാം}} picks an item out of a line: the first door, the first way. \textbf{\ml{ആദ്യം}} puts an \emph{action} first in time, and it pairs with \textbf{\ml{പിന്നെ}}, which you already have:

\begin{quote}
\textbf{\ml{ആദ്യം} \ml{ഊണ്}, \ml{പിന്നെ} \ml{ചായ}.} — \emph{Ādyaṁ ūṇŭ, pinne chāya.} \textquotedblleft{}Lunch first, then tea.\textquotedblright{}
\end{quote}

\begin{culture}
Notice what did and did not get borrowed. Malayalam's ordinals are native, and their ending is so regular that it never needed a special word for \emph{first} — the one place its sisters all reached for \textbf{\ml{മുതൽ}}. But for \emph{first} as an adverb, the language reached into Sanskrit and took \textbf{\ml{ആദ്യം}}, from \emph{ādya}, \textquotedblleft{}primary\textquotedblright{}.

So a language can have a perfect machine for one job and still borrow a word for the neighbouring one. It is worth carrying that as a habit of reading: a borrowing is evidence about a \textbf{slot}, not about the language as a whole.
\end{culture}

\begin{grammarlens}[title={What you've built}]
Two kinds of order, and the words for both:

\noindent
\begin{tabularx}{\linewidth}{@{}>{\raggedright\arraybackslash}X>{\raggedright\arraybackslash}X@{}}
\toprule
\textbf{what you are ordering} & \textbf{word} \\
\midrule
a thing in a line & \textbf{\ml{ഒന്നാം}} \emph{onnāṁ}, and the ten after it \\
an action in time & \textbf{\ml{ആദ്യം}} \emph{ādyaṁ}, with \textbf{\ml{പിന്നെ}} \emph{pinne} \\
\bottomrule
\end{tabularx}
\end{grammarlens}

\subsection*{Your turn}
Say these aloud:

\begin{itemize}
\raggedright
  \item \emph{ādyaṁ}, then \emph{ādyaṁ ūṇŭ, pinne chāya}
  \item \emph{onnāṁ vaḻi} — the first way
  \item both in one breath, and hear that they are answering different questions
\end{itemize}

\begin{tcolorbox}[breakable,colback=teal!4,colframe=teal!35!black,title={Before you move on}]
Which one picks a thing out of a line, \ml{ഒന്നാം} or \ml{ആദ്യം}? (\textbf{\ml{ഒന്നാം}}.) Which one puts an action first? (\textbf{\ml{ആദ്യം}}.) Where did \ml{ആദ്യം} come from? (\textbf{Sanskrit} \emph{ādya}.)
\end{tcolorbox}
